% Standalone solution dossier for ex:a4-gaussian-local.
\documentclass[../main.tex]{subfiles}
\begin{document}

% === SOLUTION HEADER (proof provenance) =====================================
% ledger-node: ex:a4-gaussian-local
% refines: ex:a4-gaussian-local
% bounded_by: none
% evidence_run: none
% checked_by: agent
% author: /root/a4_a5
% reviewer: /root/review_a4_a5
% dossier_editor: /root/review_a4_a5
% review: research/reviews/2026-08-21-a4-a5-independent-agent-audit.md
% date: 2026-08-21
% ============================================================================

\section*{Solution: exact fixed-covariance Gaussian calibration}

\begin{theorem}
Let $\pi=N(\mu,\Sigma)$ and
$\mathcal Q_S=\{q_m=N(\mu+m,S):m\in\R^d\}$, with $S,\Sigma\succ0$. Put
\[
 \delta_S=\tfrac12\left\{\operatorname{tr}(\Sigma^{-1}S)-d
 +\log\frac{\det\Sigma}{\det S}\right\},
\]
\[
 D_S=\operatorname{tr}\left[S+\Sigma
 -2(\Sigma^{1/2}S\Sigma^{1/2})^{1/2}\right],
 \qquad\lambda=\lambda_{\max}(\Sigma).
\]
If $S\ne\Sigma$, then for every $\rho\ge0$,
\[
 C_{\mathcal Q_S,\delta_S+\rho}(\pi)
 =\max\left\{\frac{D_S}{2\delta_S},
 \frac{D_S+2\rho\lambda}{2(\delta_S+\rho)}\right\},
 \qquad
 C_{\mathrm{mean},\mathcal Q_S,\delta_S+\rho}(\pi)
 =\lambda\frac{\rho}{\delta_S+\rho}.
\]
For every $\rho>0$, both optimizer-centred excess constants equal $\lambda$. If
$S=\Sigma$, both ordinary localized constants equal $\lambda$ for every $\rho>0$.
\end{theorem}

\begin{proof}
The Gaussian formulas give
\[
 2\KL(q_m\|\pi)=2\delta_S+m^\top\Sigma^{-1}m,
 \qquad W_2^2(q_m,\pi)=D_S+\norm m^2.
\]
Write $s=m^\top\Sigma^{-1}m$. At fixed $s$, the maximum of $\norm m^2$ is $\lambda s$,
and the localized sublevel is $0\le s\le2\rho$. Hence the transport supremum is the maximum
of $(D_S+\lambda s)/(2\delta_S+s)$ on this interval. Its derivative has constant sign, so the
maximum is at an endpoint. The mean ratio $\lambda s/(2\delta_S+s)$ is increasing, yielding
the second formula. Relative to $q_0=N(\mu,S)$, the squared transport and squared mean
displacements are both $\norm m^2$, while twice the excess KL is
$m^\top\Sigma^{-1}m$; their supremum is $\lambda$. When $S=\Sigma$, the same quotient is the
ordinary ratio.
\end{proof}

\end{document}
